\scanpage{67}
\mspage{1}
Let $K$ be a subfield of the field of complex numbers $\mathbb C$ -- interesting cases will be the one where $K=\mathbb Q$, or $K$ a finite extension of $\mathbb Q$, or $K=\widetilde{\mathbb Q}$ the maximal abelian extension of $\mathbb Q$ in $\mathbb C$ -- I shall in any case be interested above all in the cases where $K$ is of finite type over $\mathbb Q$. (The case of $\widetilde{\mathbb Q}$ is a ``limit'' case of such cases, so as not to be bothered by the roots of $1\ldots$).

We are interested in schemes \inserted{$X$} of finite type over $K$, and in the corresponding \inserted{analytic} spaces $X(\mathbb C)$. As we are interested in the ``topological'' invariants of $X,X(\mathbb C)$, we shall treat the morphisms $X'-X$ which induce universal homeomorphisms as isomorphisms -- in particular nilpotent elements do not count -- a surjective closed immersion will be treated as an isomorphism.

A convenient language is to regard $X$ as given by the analytic space (or $\overline{\mathrm m}$ merely topological) $X(\mathbb C)$, plus an additional \uline{structure} on it: the structure of an ``algebraic variety over $K$'', by means of Zariski topology, of the sheaf of $K$-polynomial functions and so on. In what follows we shall write $X$ instead of $X(\mathbb C)$.

We distinguish on $X$ the fundamental groupoid $\Pi_1(X)$ in the topological sense, and its profinite compactification $\widehat\Pi_1(X)$, which \deleted{is \unsure{in a} sense} \inserted{can} also be defined directly in terms of the $\mathbb C$-scheme $X$. We have an ``inclusion'' functor
\margin{66}
